\chapter{SPIN}\label{spin}

\hfill\break

Dopo alcuni chilometri, a distanza di sicurezza dalla fattoria di
Condon, parcheggiai e ordinai a Simon di scendere.

«Cosa?» ribatté «Qui?»

«Devo visitare Diane. Mi serve che tu prenda la torcia elettrica dal
bagagliaio e me la tenga ferma. Va bene?»

Annuì, gli occhi spalancati.

Diane non aveva detto una parola da quando avevamo lasciato il ranch.
Era rimasta sdraiata sul sedile posteriore con la testa poggiata sul
grembo di Simon, annaspando per respirare. Il suo affanno era stato il
rumore più forte dentro la macchina.

Mentre Simon aspettava in piedi reggendo la torcia, mi tolsi i vestiti
inzuppati di sangue e mi lavai il più accuratamente possibile - una
bottiglia d'acqua minerale con un po' di benzina per rimuovere lo
sporco, una seconda bottiglia per risciacquare. Poi indossai dei Levi's
puliti e una felpa presi dalla valigia e un paio di guanti di lattice
che sfilai dalla valigetta dei medicinali. Bevvi una terza bottiglia
d'acqua fino in fondo. Infine chiesi a Simon di puntare la luce su Diane
mentre le davo un'occhiata.

Era quasi cosciente ma troppo intontita per riuscire a mettere insieme
una frase con un senso logico. Era più magra che mai, quasi anoressica,
e pericolosamente febbricitante. La pressione sanguigna e il polso erano
elevati, e quando le auscultai il petto sentii i suoi polmoni emettere
un rumore simile a quello di un bambino che succhia un frappè da una
cannuccia stretta.

Riuscii a farle inghiottire un po' d'acqua e un'aspirina. Poi strappai
il sigillo di una siringa ipodermica sterile.

«Cos'è quello?» chiese Simon.

«Un antibiotico generico.» Le passai un tampone sul braccio e con
qualche difficoltà trovai una vena. «Ne servirà una anche a te.» E a me.
Il sangue della giovenca era senza dubbio carico di batteri vivi di
sindrome da deperimento cardiovascolare.

«Questo la curerà?»

«No, Simon, temo di no. Un mese fa avrebbe potuto. Ora non più. Le
servono cure mediche.»

«Tu sei un medico.»

«Sarò anche un medico, ma non sono un ospedale.»

«Allora forse potremmo portarla a Phoenix.»

Ci avevo pensato. Tutto ciò che avevo imparato durante gli sfarfallii
era che un ospedale urbano sarebbe stato sommerso di lavoro nella
migliore delle ipotesi, un mucchio di rovine fumanti nella peggiore. Ma
forse no.

Afferrai il telefono e controllai la rubrica alla ricerca di un numero
che avevo quasi dimenticato di avere.

Simon si incuriosì. «A chi stai telefonando?»

«A qualcuno che conoscevo un tempo.»

Si chiamava Colin Hinz ed eravamo stati coinquilini a Stony Brook. Ci
eravamo tenuti in contatto per un po'. L'ultima volta che l'avevo
sentito lavorava al Saint Joseph a Phoenix. Valeva la pena di tentare -
ora, prima che il Sole sorgesse e cancellasse le telecomunicazioni per
un altro giorno.

Chiamai il suo numero privato. Il telefono squillò a lungo ma alla fine
rispose: «Sarà meglio che sia importante.»

Mi identificai e gli dissi che mi trovavo a circa un'ora dalla città con
un paziente che aveva bisogno di cure urgenti - e che era qualcuno a me
molto caro.

Colin sospirò. «Non so cosa dirti, Tyler. Il Saint Joe è attivo, e mi
hanno detto che la Mayo Clinic a Scottsdale è aperta, ma siamo entrambi
a corto di personale. Dagli altri ospedali arrivano informazioni
contrastanti. Ma non riuscirai a ricevere cure rapide da nessuna parte,
di sicuro non qui. Abbiamo gente ammassata fuori dalle porte - ferite da
arma da fuoco, tentati suicidi, incidenti d'auto, attacchi di cuore, di
tutto un po'. E poliziotti all'ingresso per evitare che venga assaltato
il reparto delle emergenze. In che condizioni è la tua paziente?»

Gli dissi che Diane era nella fase terminale della Sindrome da
Deperimento Cardiovascolare e forse presto avrebbe avuto bisogno di
essere intubata per respirare.

«Dove cazzo ha preso la SDC? No, lascia perdere\ldots{} non importa.
Onestamente, ti aiuterei se potessi, ma le nostre infermiere stanno
facendo il triage nel parcheggio da tutta la notte e non posso
prometterti che daranno la priorità alla tua paziente, nemmeno se ci
metto una buona parola io. In realtà, è quasi certo che non verrà
nemmeno valutata da un medico prima di altre ventiquattr'ore. Sempre che
qualcuno di noi viva tanto a lungo.»

«Io sono un medico, ricordi? Mi serve solo un po' d'attrezzatura per
tenerla in vita. Soluzione salina, un kit per l'intubazione,
ossigeno\ldots»

«Non vorrei sembrarti insensibile, ma qui stiamo sguazzando nel
sangue\ldots{} prova a chiederti se vale davvero la pena di assistere un
caso terminale di sindrome da deperimento cardiovascolare, con quello
che sta succedendo. Se hai ciò che ti serve per metterla a proprio
agio\ldots»

«Non voglio metterla a proprio agio. Voglio salvarle la vita.»

«D'accordo\ldots{} ma quella che mi hai descritto è una situazione
terminale, a meno che non abbia frainteso.» In sottofondo sentivo altre
voci che richiamavano la sua attenzione, un generico rantolo di miseria
umana.

«Devo portarla in un posto,» dissi, «e devo farcela arrivare viva. Più
che un letto mi serve il materiale sanitario.»

«Non abbiamo scorte. Dimmi se c'è altro che posso fare, altrimenti, mi
dispiace, ma ho del lavoro da fare.»

Pensai freneticamente. Poi risposi: «Va bene, ma il materiale\ldots{}
dimmi solo dove posso trovare della soluzione salina, Colin, non chiedo
altro.»

«Be'\ldots»

«Be', cosa?»

«Be'\ldots{} non dovrei dirtelo, ma il Saint Joe ha un accordo con la
città per le emergenze. C'è un distributore medico chiamato Novaprod a
nord della città.» Mi diede un indirizzo e alcune semplici istruzioni.
«Le autorità hanno messo a proteggerlo un'unità della Guardia Nazionale.
È quella la nostra fonte primaria di medicinali e macchinari.»

«Mi faranno entrare?»

«Se li chiamo e li avviso che stai arrivando, e se hai una carta
d'identità da mostrare.»

«Fallo per me, Colin. Ti prego.»

«Lo farò se riesco ad avere segnale. I telefoni non sono affidabili.»

«Se c'è qualcosa che posso fare per sdebitarmi\ldots»

«Forse c'è. Lavoravi nel campo aerospaziale, vero? Perielio?»

«Non di recente, ma sì.»

«Mi puoi dire quanto durerà tutto questo?» Pose la domanda quasi
bisbigliando, e di colpo riuscii a sentire la fatica nella sua voce, la
paura inconfessabile. «Voglio dire, in un modo o nell'altro?»

Mi scusai, gli risposi che proprio non ne avevo idea - e dubitavo che
qualcuno alla Perielio ne sapesse più di me.

Sospirò. «Va bene» disse. «È solo che mi secca l'idea che dobbiamo
affrontare tutto questo per poi, magari, bruciare in un paio di giorni
senza sapere cosa sia successo.»

«Vorrei poterti rispondere.»

Qualcuno dall'altra parte della linea cominciò a chiamarlo. «Anch'io
vorrei un sacco di cose» mormorò. «Devo andare, Tyler.»

Lo ringraziai di nuovo e chiusi la conversazione.

Mancavano ancora poche ore all'alba.

Simon era rimasto in piedi a pochi metri dalla macchina, fissando il
cielo stellato e facendo finta di non sentire. Gli feci cenno di
avvicinarsi. «Dobbiamo andare» gli dissi.

Annuì docilmente. «Hai trovato aiuto per Diane?»

«Più o meno.»

Accettò la risposta senza chiedere altri dettagli. Ma prima di chinarsi
per entrare in macchina, mi afferrò per una manica.
«\emph{Laggiù\ldots{} cosa pensi che sia, Tyler}?» chiese.

Stava indicando l'orizzonte occidentale, dove nel cielo notturno una
linea argentata si incurvava dolce in un arco di cinque gradi. Sembrava
che qualcuno avesse inciso nell'oscurità un'enorme e delicata lettera C.

«Forse è una scia di condensazione» ipotizzai. «Un jet militare.»

«Di notte? No, non di notte.»

«Allora non so cosa sia, Simon. Forza, entra\ldots{} non c'è tempo da
perdere.»

ooo

Ci mettemmo meno tempo di quanto avessi previsto. Arrivammo al magazzino
di forniture mediche, un lotto numerato in una tetra zona industriale,
in anticipo sul sorgere del Sole. All'ingresso presentai il documento
d'identità a una guardia che sembrava nervosa; lui mi consegnò a un
collega e a un impiegato civile che mi accompagnò per i corridoi tra gli
scaffali. Trovai ciò che mi serviva e una terza guardia mi aiutò a
portare tutto in macchina, ma indietreggiò con uno scatto quando vide
Diane che rantolava sul sedile posteriore. «Buona fortuna» mormorò con
voce leggermente tremula.

Mi presi del tempo per allestire una flebo, agganciando la sacca alla
meglio sul poggiamano superiore dell'auto, e mostrai a Simon come tenere
sotto controllo il flusso e assicurarsi che Diane non staccasse il
tubicino nel sonno.

(Non si svegliò nemmeno quando le misi l'ago nel braccio.)

Simon aspettò che fossimo di nuovo in strada poi mi chiese: «Sta
morendo?»

Strinsi il volante più forte. «Non se posso evitarlo.»

«Dove la stiamo portando?»

«La stiamo portando a casa.»

«Cosa? E attraversiamo tutto il paese? A casa di Carol e E.D.?»

«Esatto.»

«Perché proprio là?»

«Perché laggiù posso aiutarla.»

«È un viaggio lungo. Voglio dire, visto come stanno le cose.»

«Sì. Potrebbe essere un viaggio lungo.»

Gettai lo sguardo sul sedile posteriore. Le accarezzava la testa,
delicatamente.

Lei aveva i capelli flosci e arruffati dal sudore. Lui aveva le mani
pallide dove aveva lavato via il sangue.

«Non merito di stare insieme a lei» disse. «Lo so che è colpa mia. Avrei
potuto lasciare il ranch insieme a Teddy. Avrei potuto trovare aiuto.»

``Sì'', pensai, ``avresti potuto''.

«Ma credevo fortemente in ciò che stavamo facendo. Forse tu non lo
capisci. Ma non era soltanto il vitello rosso, Tyler. Ero certo che
saremmo stati assunti in cielo. Che alla fine saremmo stati
ricompensati.»

«Ricompensati per cosa?»

«Fede. Perseveranza. Perché dalla prima volta che la vidi ebbi la forte
sensazione che saremmo stati partecipi di qualcosa di straordinario,
anche se non mi era chiaro di cosa si trattasse. Che un giorno ci
saremmo trovati di fronte al trono di Dio - proprio così. ``Questa
generazione non passerà prima che queste cose siano avvenute''. La
\emph{nostra} generazione, anche se all'inizio abbiamo imboccato la via
sbagliata. Lo ammetto, a quei raduni del Nuovo Regno sono successe cose
di cui ora mi vergogno. Ubriachezza, libidine, bugie. Ce le siamo
lasciate alle spalle quelle cose, e abbiamo fatto bene; ma quando ci
siamo chiamati fuori da quelli che, seppure in modo imperfetto,
provavano a costruire il chiliasmo, ci è sembrato che il mondo si fosse
un po' rimpicciolito. Come se avessimo perso una famiglia. Allora mi son
detto: bene, se cerchi il sentiero più pulito e semplice, quello ti
porterà nella giusta direzione. ``\emph{Con la vostra perseveranza
	guadagnerete le vostre anime}''.»

«Il Jordan Tabernacle» dissi.

«È facile profetizzare contro lo Spin. Segni sul Sole e sulla Luna e
sulle stelle, lo dice il Vangelo di Luca. Bene, eccoci qui. Le forze del
cielo sono state scosse. Ma non è\ldots{} non è che\ldots»

Sembrò aver perso il filo.

«Come sta respirando Diane, là dietro?» Ma non avevo davvero bisogno di
chiederlo. Potevo sentirne ogni respiro, faticoso ma regolare. Volevo
solo distrarlo.

«Non sta soffrendo» rispose Simon. Poi aggiunse: «Ti prego, Tyler.
Fermati e fammi scendere.»

Stavamo andando a est. Sorprendentemente c'era poco traffico
sull'interstatale. Colin Hinz mi aveva avvertito dell'ingorgo attorno
all'aeroporto Sky Harbor, ma l'avevamo evitato. Là fuori avevamo
incrociato solo poche auto con qualcuno a bordo. I veicoli abbandonati
sul ciglio della strada erano molti di più.

«Non è una buona idea» lo redarguii.

Guardai nello specchietto e vidi Simon che si asciugava le lacrime con
le nocche. In quel momento sembrava vulnerabile e confuso come un bimbo
di dieci anni a un funerale.

«Ho avuto solo due punti di riferimento nella mia vita» disse. «Dio e
Diane. E li ho traditi entrambi. Ho aspettato troppo. Sei gentile a
negarlo, ma sta morendo.»

«Non necessariamente.»

«Non voglio stare con lei sapendo che avrei potuto prevenire tutto
questo. Voglio solo morire al più presto, nel deserto. Lo penso davvero,
Tyler. Voglio uscire.»

Il cielo si stava schiarendo di nuovo, un minaccioso bagliore violetto
più simile a un arco in una lampada fluorescente guasta che a qualcosa
di salubre o naturale.

«Non m'interessa» puntualizzai.

Simon mi lanciò uno sguardo stupefatto. «Cosa?»

«Non m'interessa come ti senti. Il motivo per cui devi starle accanto è
che dobbiamo ancora fare molta strada e io non posso guidare e prendermi
cura di lei contemporaneamente. E prima o poi dovrò pure dormire. Se una
volta ogni tanto guidi tu, non dovremo fermarci se non per mangiare e
fare benzina,» sempre che ne trovassimo, «se te ne vai, la durata del
viaggio raddoppia.»

«E quindi? Cosa importa?»

«Forse non starà morendo, Simon, ma è davvero malata come pensi, e
morirà se non riceverà delle cure. E l'unico aiuto che conosco si trova
a circa tremila chilometri da qui.»

«Il cielo e la terra sono al trapasso. Moriremo tutti.»

«Non posso parlare per il cielo e la terra. Ma finché posso scegliere,
mi rifiuto di farla morire.»

«Ti invidio» mormorò Simon.

«Cosa? Che cosa potresti mai invidiarmi?»

«La tua fede» rispose.

ooo

Un certo tipo di ottimismo era ancora possibile, ma solo di notte. Di
giorno appassiva.

Guidai verso est di Flagstaff, nell'Hiroshima del Sole nascente. Avevo
smesso di preoccuparmi che la luce stessa potesse uccidermi, sebbene con
tutta probabilità non mi facesse bene. Il fatto che ognuno di noi fosse
sopravvissuto il primo giorno era un mistero - un miracolo, come lo
avrebbe definito Simon. Mi spronò a un senso pratico alquanto rozzo:
sfilai dal cruscotto un paio di occhiali da sole e provai a tenere lo
sguardo fisso sulla strada invece che sull'emisfero di fuoco arancione
che saliva all'orizzonte.

Il giorno si fece più caldo. Così come l'interno della macchina,
nonostante il condizionatore sovraccarico di lavoro. (Facevo andare
l'aria al massimo nel tentativo di tenere sotto controllo la temperatura
corporea di Diane.) Da qualche parte tra Albuquerque e Tucumcari fui
travolto da un'enorme ondata di stanchezza. Mi si chiudevano le palpebre
e quasi andai a sbattere contro un segnale stradale di progressiva
chilometrica. A quel punto accostai e spensi il motore. Ordinai a Simon
di riempire il serbatoio con le taniche e di prepararsi a sostituirmi
alla guida. Annuì riluttante.

Eravamo in anticipo sulla tabella di marcia. Il traffico era stato
leggero, quasi inesistente, forse perché le persone avevano paura di
stare da sole in strada. Mentre Simon metteva la benzina nell'auto, gli
chiesi: «Cosa hai portato da mangiare?»

«Solo quello che sono riuscito ad arraffare dalla cucina. Dovevo
sbrigarmi. Dai un'occhiata.»

Trovai una scatola di cartone tra le taniche ammaccate, le forniture
mediche imballate e le bottiglie d'acqua sparse nel bagagliaio.
Conteneva tre scatole di Cheerios, due lattine di carne in scatola e una
bottiglia di Diet Pepsi. «Cristo, Simon.»

Fece una smorfia nel sentire quella che, dovetti ricordare a me stesso,
per lui era una bestemmia. «È tutto ciò che sono riuscito a trovare.»

E nemmeno una scodella o dei cucchiai. Ma ero affamato e in carenza di
sonno. Spiegai a Simon che dovevamo lasciare che il motore si
raffreddasse, abbassammo i finestrini e ci mettemmo a sedere all'ombra
della macchina, una brezza decisa proveniente dal deserto, il Sole
sospeso in cielo come un mezzogiorno sulla superficie di Mercurio.
Usammo i fondi tagliati dalle bottiglie di plastica vuote come tazze di
fortuna e mangiammo i Cheerios ammorbiditi con acqua tiepida. Sembrava
mucillagine e aveva lo stesso sapore.

Ragguagliai Simon sulla successiva tappa del viaggio, gli ricordai di
accendere l'aria condizionata una volta in movimento, gli dissi di
svegliarmi in caso gli sembrasse di vedere problemi in strada.

Poi mi occupai di Diane. La flebo e gli antibiotici parevano averla
temprata, ma solo in misura marginale. Aprì gli occhi e dopo che la
aiutai a bere un sorso d'acqua disse: «Tyler.» Accettò qualche cucchiaio
di Cheerios ma subito dopo girò la testa dall'altra parte. Aveva le
guance scavate, gli occhi spenti e assenti.

«Porta pazienza» la esortai. «Ancora per un po', Diane.» Le aggiustai la
flebo. La aiutai a tirarsi su, a gambe divaricate fuori dalla macchina,
mentre espelleva dell'urina brunastra. Poi la pulii con una spugnetta e
le sostituii le mutande sporche con un paio di slip di cotone presi
dalla mia valigia.

Quando fu di nuovo comoda, riempii lo spazio angusto tra i sedili
anteriori e posteriori con una coperta, in modo da potermi distendere
senza spostarla. Simon aveva fatto solo un rapido sonnellino durante la
prima parte del viaggio e doveva essere stanco quanto me\ldots{} ma non
era stato colpito dal calcio di un fucile. Il punto in cui fratello
Aaron mi aveva percosso era gonfio e, quando ci avvicinavo le dita, il
dolore era sordo come una campana.

Simon osservava la scena a un paio di metri di distanza, con espressione
imbronciata, o forse gelosa. Quando lo chiamai, esitò e fissò bramoso il
deserto salino nel cuore profondo del nulla assoluto.

Poi si diresse a grandi passi verso la macchina, lo sguardo basso, e si
mise al volante.

Mi incastrai nella mia nicchia dietro il sedile anteriore. Diane pareva
incosciente, ma prima di addormentarmi la sentii premere la sua mano
sulla mia.

ooo

Quando mi risvegliai, era di nuovo notte e Simon aveva accostato per
scambiarci di posto.

Saltai fuori dall'auto e mi stiracchiai. La testa continuava a pulsare e
sentivo la spina dorsale contratta come quella di un anziano, ma ero più
sveglio di Simon, che strisciò dietro e si addormentò all'istante.

Non sapevo dove fossimo, ma solo che ci trovavamo sull'Interstatale 40
in direzione est e lì la terra era meno arida. Su entrambi i lati della
strada e sotto la luce cremisi della Luna si stendevano campi irrigati.
Mi assicurai che Diane fosse comoda e respirasse senza difficoltà e
lasciai aperti gli sportelli anteriori e posteriori per un paio di
minuti per mandare via la puzza, un insalubre miasma di infermeria
sovrapposto a odore di sangue e benzina. Poi mi misi al posto del
guidatore.

Le stelle sopra la strada erano tristemente poche e impossibili da
riconoscere. Mi chiedevo di Marte. Era ancora sotto una membrana dello
Spin o era stato liberato come la Terra? Ma non sapevo in quale parte
del cielo cercare e dubitavo di riuscire a individuarlo, anche se
l'avessi visto. Notai - non potevo farne a meno - la linea argentata che
Simon mi aveva indicato in Arizona, quella che avevo scambiato per una
scia di condensazione. Stanotte era ancora più evidente. Si era mossa
dall'orizzonte occidentale fino quasi allo zenit, e la leggera curvatura
era diventata ovale, una lettera O appiattita.

Il cielo che osservavo era di tre miliardi di anni più vecchio di quello
che avevo visto l'ultima volta dal prato della Grande Casa. Presumevo
che nascondesse ogni genere di mistero.

Una volta in marcia armeggiai con la radio del cruscotto, che la notte
precedente era rimasta muta. Non trovai alcuna trasmissione digitale ma
finalmente rintracciai una stazione locale sulla banda FM - quel genere
di stazione di paese solitamente dedicata alla musica country e alla
cristianità, che però quella notte trasmetteva solo informazioni.
Imparai un sacco di cose, prima che il segnale svanisse del tutto in un
rumore indistinto.

Per prima cosa appresi che eravamo stati saggi a evitare le grandi
città. Le città principali erano zone disastrate - non per saccheggi e
violenze (ce n'erano stati sorprendentemente pochi) ma a causa del
crollo catastrofico delle infrastrutture. Il sorgere del Sole rosso
assomigliava tanto alla morte della Terra prevista da lungo tempo che la
maggior parte delle persone erano rimaste a casa per morire con le
proprie famiglie, lasciando i centri urbani con un numero minimo di
poliziotti e pompieri e gli ospedali con personale assai ridotto. La
minoranza di persone che cercarono di uccidersi a colpi di pistola, o
che si stordirono con dosi smodate di alcool, cocaina, ossicodone o
anfetamine, furono involontariamente la causa dei problemi più
imminenti: lasciarono le stufe a gas in funzione, persero conoscenza
guidando o lasciarono cadere le sigarette mentre morivano. Quando la
moquette cominciò a incendiarsi o le tende presero fuoco nessuno chiamò
il 911, e in molti casi non ci sarebbe stato nessuno dall'altra parte a
rispondere. Gli incendi delle case presto diventarono incendi di
quartieri.

Quattro enormi spire di fumo si erano sollevate da Oklahoma City, disse
il giornalista, e secondo le segnalazioni telefoniche gran parte
dell'area sud di Chicago era già ridotta in cenere. Ogni grande città di
questo paese - tutte quelle che erano state interpellate - segnalava
almeno uno o due incendi su vasta scala fuori controllo.

Ma la situazione stava migliorando, non peggiorando. Cominciava a
sembrare possibile che la specie umana potesse sopravvivere almeno per
un altro po', e di conseguenza molti addetti al primo soccorso e ai
servizi essenziali erano tornati al lavoro. (Il lato negativo stava nel
fatto che la popolazione aveva cominciato a preoccuparsi di quanto
sarebbero durate le provviste e a saccheggiare i negozi di alimentari.)
Chiunque non dovesse garantire un servizio essenziale veniva invitato a
stare lontano dalle strade - il messaggio era stato trasmesso prima
dell'alba dal sistema di trasmissione d'emergenza e dalle radio e TV
ancora funzionanti, e quella notte venne spesso replicato. Ecco perché
il traffico sull'interstatale era stato piuttosto limitato. Avevo visto
qualche pattuglia di militari e polizia ma nessuno di loro ci aveva
ostacolato, presumibilmente grazie alle targhette sulla macchina - dopo
il primo episodio di sfarfallio, la California e la maggior parte degli
altri stati avevano iniziato a distribuire etichette per i medici
dell'EMS, il servizio medico di emergenza, da apporre sulla macchina.

La vigilanza della polizia era sporadica. I militari regolari rimanevano
quasi gli stessi a parte qualche diserzione, ma i reparti della Riserva
e della Guardia Nazionale erano ridotti ai minimi termini e non
riuscivano a sostituirsi alle autorità locali. Era sporadica anche la
corrente elettrica; la maggior parte delle centrali era a corto di
personale e funzionava a malapena, e lungo la rete erano cominciati i
blackout. Girava voce che gli impianti nucleari di San Onofre in
California e di Pickering in Canada fossero arrivati a un passo dalla
fusione del nocciolo, sebbene la cosa non fosse stata confermata.

L'annunciatore poi elencò una lista di depositi locali di prodotti
alimentari istituiti ad hoc, ospedali ancora aperti al pubblico (con
stime sui tempi di attesa per il triage) e consigli utili di primo
soccorso da fare a casa. Lesse anche un documento consultivo
dell'Osservatorio Meteorologico che metteva in guardia dall'esposizione
prolungata al Sole. La luce solare non era immediatamente fatale, ma
livelli eccessivi di ultravioletti potevano causare ``problemi a lungo
termine'', dissero, cosa che risultò triste e buffa allo stesso tempo.

ooo

Fino all'alba intercettai qua e là qualche altra trasmissione, ma il
sorgere del Sole mise fine ai programmi, sostituiti da un fruscio
indistinto.

Il giorno in arrivo era nuvoloso. Quindi, non avrei avuto il problema di
guidare con la luce del Sole negli occhi; ma perfino quell'alba sopita
era strana e desolante. L'intera metà orientale del cielo divenne una
brodaglia di luce rossa, a suo modo ipnotica come le braci di un falò
che si stanno spegnendo. Di tanto in tanto le nuvole si aprivano e dita
di luce ambrata sondavano il terreno. Ma a mezzogiorno la nuvolosità si
intensificò e nel giro di un'ora iniziò a piovere - una pioggia calda e
senza vita che ricopriva l'autostrada e rispecchiava i colori malaticci
del cielo.

Quella mattina avevo svuotato l'ultima tanica di benzina nel serbatoio,
e da qualche parte tra Cairo e Lexington la lancetta del carburante
aveva cominciato a scendere in maniera preoccupante. Svegliai Simon, lo
misi al corrente del problema e lo informai che sarei entrato nella
prima stazione di servizio che avremmo incontrato\ldots{} e in ognuna di
quelle successive finché non ne avessimo trovata una che ci vendesse del
carburante.

La prima stazione in cui ci imbattemmo si rivelò essere un piccolo
distributore a quattro pompe a conduzione familiare di una catena di
autogrill, a quattrocento metri dall'autostrada. Il negozio era buio e
le pompe probabilmente esaurite, ma mi avvicinai lo stesso, scesi
dall'auto e tolsi l'erogatore dal suo gancio.

Un uomo col cappellino dei Bengals e un fucile imbracciato arrivò da un
lato dell'edificio e disse: «Non funziona.»

Rimisi a posto il beccuccio, lentamente. «Manca la corrente?»

«Esatto.»

«Nessuna soluzione alternativa?»

Scrollò le spalle e si avvicinò. Simon stava per scendere dalla macchina
ma gli feci segno di tornare dentro. L'uomo col cappellino dei Bengals -
sui trent'anni e sovrappeso di una decina di chili - lanciò un'occhiata
alla flebo allestita nel sedile posteriore. Poi strizzò gli occhi per
leggere la targa. Era una targa della California, cosa che probabilmente
non mi fece guadagnare punti benevolenza, ma l'adesivo EMS spiccava
chiaramente. «Sei un dottore?»

«Tyler Dupree,» risposi, «e sì, sono un medico.»

«Scusami se non ti stringo la mano. Quella là è tua moglie?»

Dissi di sì, perché era più facile che dare una spiegazione. Simon mi
lanciò un'occhiataccia ma non mi contraddisse.

«Hai un documento che provi che sei un medico? Perché, senza offesa, ci
sono stati parecchi furti d'auto in questi ultimi giorni.»

Tirai fuori il portafogli e lo lanciai ai suoi piedi. Lo raccolse e
frugò nello scomparto delle carte. Poi pescò un paio d'occhiali dalla
tasca della camicia e continuò a frugare. Alla fine me lo restituì e mi
porse la mano. «Mi scusi, dottor Dupree. Sono Chuck Bernelli. Se l'unica
cosa che vi serve è la benzina vi accendo le pompe. Se vi serve
dell'altro ci metto un minuto ad aprire il negozio.»

«Mi serve la benzina. Mi farebbero comodo anche delle provviste, ma non
ho molto denaro liquido.»

«Chi se ne frega dei soldi. Siamo chiusi per i criminali e gli ubriachi,
che ormai sembrano essere gli unici ad andare in giro, ma per i militari
e le pattuglie autostradali siamo sempre aperti. E anche per i medici.
Almeno finché c'è carburante nelle pompe. Spero che tua moglie non sia
in condizioni troppo gravi.»

«Non se riesco ad arrivare dove sono diretto.»

«Lexington V.A.? Samaritan?»

«Un po' più lontano. Ha bisogno di cure speciali.»

Guardò ancora verso la macchina. Simon aveva abbassato i finestrini per
far entrare un po' d'aria fresca. La pioggia impastata con la polvere
dell'auto colava formando pozzanghere d'asfalto untuoso. Bernelli vide
per un attimo Diane che si girava e cominciava a tossire nel sonno. Si
incupì.

«Allora metto in funzione le pompe» disse. «Vorrai proseguire il viaggio
al più presto.»

Prima che ripartissimo, ci mise insieme un po' di generi alimentari,
qualche lattina di zuppa, una confezione di cracker salati e un
apriscatole in un imballo di plastica trasparente. Ma non volle
avvicinarsi alla macchina.

ooo

La tosse ciclica, straziante, è un comune sintomo di sindrome da
deperimento cardiovascolare. Il batterio è quasi astuto nel modo in cui
preserva le sue vittime, preferendo non affogarle in una catastrofica
polmonite, nonostante sia il mezzo con cui alla fine le uccide - quello,
o una insufficienza cardiaca completa. Dal grossista fuori Flagstaff
avevo preso una bombola d'ossigeno, una valvola di sfiato e una
maschera, e quando la tosse di Diane cominciò a interferire con la
respirazione - era sull'orlo del panico, roteava gli occhi e stava
affogando nel suo stesso muco - le ripulii al meglio le vie respiratorie
e le tenni la maschera sulla bocca e sul naso mentre Simon guidava.

Alla fine si calmò, assunse un colorito migliore e riuscì a dormire di
nuovo. Sedetti accanto a lei mentre riposava, e la sua testa
febbricitante si annidò sulla mia spalla. La pioggia era diventata un
inarrestabile acquazzone che ci rallentava. Grandi pennacchi d'acqua si
sollevavano dietro la macchina, come la coda di un gallo, ogni volta che
entravamo in un affossamento della strada. Verso sera la luce svanì in
carboni ardenti sull'orizzonte occidentale.

Non si sentiva alcun suono tranne il battito della pioggia sul tetto
dell'auto e mi piacque ascoltarlo fino a quando Simon si schiarì la voce
e mi chiese: «Tu sei ateo, Tyler?»

«Scusa?»

«Non vorrei essere scortese, ma mi domandavo: ti consideri ateo?»

Non sapevo bene come rispondere a quella domanda. Simon era stato utile
- era stato prezioso - a farci arrivare così lontano. Ma era anche uno
che aveva agganciato il suo vagone dell'intelletto a una frangia
estremista di svitati dispensazionalisti la cui unica discussione sulla
fine del mondo verteva sul fatto che avesse deluso le loro dettagliate
aspettative. Non volevo offenderlo perché avevo ancora bisogno di
lui\ldots{} \emph{Diane} aveva ancora bisogno di lui.

Quindi risposi: «Che importanza ha come mi considero?»

«Solo per curiosità.»

«Bene, non lo so. Suppongo che la mia risposta sia questa. Non pretendo
di sapere se Dio esiste o perché ha fatto l'universo a spirale per farlo
girare in quel modo. Scusa, Simon. Questo è il meglio che posso dare sul
fronte teologico.»

Rimase in silenzio per qualche altro chilometro. Poi aggiunse: «Forse è
quello che intendeva Diane.»

«Quello che intendeva su cosa?»

«Quando ne abbiamo parlato. Cosa che non abbiamo fatto di recente, ora
che ci penso. Non eravamo d'accordo sul Pastore Dan e sul Jordan
Tabernacle nemmeno prima dello scisma. Credevo che fosse troppo cinica.
Lei disse che ero facilmente impressionabile, troppo. Forse è così. Il
Pastore Dan aveva il dono di guardare nelle Scritture e trovare la
conoscenza in ogni pagina - una conoscenza solida come una casa, travi e
pilastri di conoscenza. È veramente un dono. Io non ci riesco. Per
quanto mi sforzi, non riesco ad aprire la Bibbia e trovarci un senso
immediato.»

«Forse non è compito tuo.»

«Ma lo volevo. Volevo essere come il Pastore Dan: intelligente e, sai,
sempre su un terreno solido. Diane diceva che Dan Condon aveva fatto un
patto col diavolo, che aveva barattato l'umiltà con la sicurezza. Forse
è questo che mi è mancato. Forse è questo che lei ha visto in te, il
motivo per cui ti è rimasta legata per tutti questi anni\ldots{} la tua
umiltà.»

«Simon, io\ldots»

«Non c'è niente di cui scusarsi o che tu possa dire per farmi sentire
meglio. Lo so che ti chiamava quando pensava che dormissi o quando ero
fuori casa. So di essere stato fortunato ad averla, finché è durata.» Mi
guardò di nuovo. «Potresti farmi un favore? Vorrei che le dicessi che mi
dispiace di non essermi preso maggiore cura di lei quando si è
ammalata.»

«Potrai dirglielo tu stesso.»

Annuì pensoso e continuò a guidare sotto la pioggia. Gli chiesi di
cercare qualche informazione utile alla radio, ora che era di nuovo
buio. Sarei voluto rimanere sveglio e ascoltare; ma la testa mi pulsava
e ci vedevo doppio, e dopo un po' mi sembrò più facile chiudere
semplicemente gli occhi e dormire.

ooo

Dormii profondamente e a lungo, e i chilometri passarono sotto le ruote
dell'automobile.

Quando mi svegliai era un'altra mattinata piovosa. Eravamo parcheggiati
in un'area di sosta (a ovest di Manassas, lo capii più tardi) e una
donna con un ombrello nero rotto stava picchiettando sul finestrino.

Battei gli occhi e aprii la portiera. Lei fece un passo indietro,
lanciando uno sguardo attento a Diane. «L'uomo ha detto di dirvi di non
aspettare.»

«Mi scusi?»

«Ha detto di dirvi addio e di non aspettarlo.»

Simon non era sul sedile davanti. Non era nemmeno visibile tra i bidoni
della spazzatura, i tavolini da picnic fradici e le latrine poco
allettanti nelle immediate vicinanze. C'erano poche altre macchine
parcheggiate, gran parte delle quali con il motore al minimo mentre i
proprietari facevano una capatina ai bagni. Mi guardai attorno: degli
alberi, un parco, la vista collinare di una cittadina industriale
infradiciata dalla pioggia sotto un cielo ardente. «Un tipo magro e
biondo? Maglietta sporca?»

«Esatto. Proprio lui. Ha detto che non voleva che dormiste troppo a
lungo. Poi è sparito.»

«A piedi?»

«Sì. Giù verso il fiume, non lungo la strada.» Guardò di nuovo Diane che
respirava a malapena e rumorosamente. «Voi due state bene?»

«No. Ma non dobbiamo andare troppo lontano. Grazie per averlo chiesto.
Ha detto nient'altro?»

«Sì. Mi ha detto di dirvi ``che Dio vi benedica'' e che da qui in avanti
avrebbe trovato da solo la sua strada.»

Mi occupai dei bisogni di Diane. Diedi un'ultima occhiata in giro per il
parcheggio piovoso. Poi tornai in strada.

ooo

Mi dovetti fermare svariate volte per sistemare la flebo a Diane o farle
prendere qualche boccata d'ossigeno. Non apriva più gli occhi - non era
semplicemente addormentata, aveva perso conoscenza. Mi rifiutavo di
pensare a cosa volesse significare.

Il viaggio andava a rilento, la pioggia era incessante e ovunque c'erano
tracce del caos degli ultimi due giorni. Superai decine di macchine
sfasciate o abbandonate e incendiate ai lati della strada, alcune ancora
fumanti. Certi tratti erano stati chiusi al traffico dei civili e
riservati ai veicoli militari o d'emergenza. Dovetti tornare indietro
dai posti di blocco un paio di volte. Il calore diurno rendeva l'aria
umida quasi insopportabile e, sebbene nel pomeriggio si fosse alzato un
vento impetuoso, non aveva portato sollievo.

Ma per fortuna Simon ci aveva abbandonato vicino alla nostra
destinazione, e riuscii ad arrivare alla Grande Casa quando c'era ancora
un po' di luce nel cielo.

Il vento era peggiorato, con raffiche quasi da record, e il lungo viale
d'accesso dei Lawton era disseminato di rami strappati dai pini
circostanti. La casa stessa era scura, o così sembrava, immersa nel
tramonto ambrato.

Lasciai Diane in macchina ai piedi delle scale e bussai alla porta. E
aspettai. E bussai ancora. Alla fine la porta si schiuse e Carol Lawton
fece capolino.

Riuscivo a malapena a scorgerne il viso attraverso quella fessura: un
occhio celeste, un lembo di guancia rugosa. Ma lei mi riconobbe.

«Tyler Dupree!» esclamò. «Sei da solo?»

La porta si spalancò.

«No,» risposi, «c'è Diane con me. E mi servirà aiuto per portarla
dentro.»

Carol uscì nel grande portico anteriore e guardò verso la macchina.
Quando vide Diane, il suo corpo esile si irrigidì; alzò le spalle e
sospirò.

«Mio Dio» sussurrò. «\emph{Entrambi} i miei figli sono venuti a morire a
casa?»
